% ============================================================
%  Chapter 9 — Robot Friends: Companions and Pets
%  Robot World: Meet the Machines That Help Us
% ============================================================

\chapter{Robot Friends \textemdash{} Companions and Pets}

\chapteropening{%
  Can a robot be your friend? Some robots are built not to work,
  but to keep us company, make us smile, and even help us feel less alone.
}
\chapterbadge{Cora: Learning to Be a Friend}

% --- Introduction ---
\section*{More Than Metal}

\character{Cora}{I wave when you come home, tilt my head when you talk,
  and learn your favourite jokes. I am not alive, but I try to make
  your day a little brighter.}

Companion robots are designed to interact with people socially. They
recognise faces, respond to voices, and express emotions through lights,
sounds, and movement.

\sectionrule

% --- Famous companion robots ---
\section*{Famous Companion Robots}

\begin{itemize}
  \item \textbf{\term{AIBO}} \textemdash{} Sony's robot dog. It wags its tail,
    learns tricks, and develops a personality over time.
  \item \textbf{\term{Paro}} \textemdash{} a fluffy robotic seal used in care homes.
    Stroking it lowers stress and brings comfort to elderly residents.
  \item \textbf{\term{Pepper}} \textemdash{} a humanoid robot that reads facial
    expressions and adjusts its behaviour accordingly.
  \item \textbf{Miko} \textemdash{} a small educational companion that helps
    children learn through conversation and games.
  \item \textbf{Vector} \textemdash{} a palm-sized robot that explores your desk,
    answers questions, and reacts to touch.
\end{itemize}

\begin{funfact}
  When Sony stopped repairing AIBO robots, owners in Japan held
  funerals for their robot pets at a Buddhist temple.
  That is how attached people became!
\end{funfact}

\sectionrule

% --- How they work ---
\section*{How Companion Robots Connect}

Companion robots use:
\begin{itemize}
  \item \textbf{\term{Face recognition}} \textemdash{} cameras identify who is talking.
  \item \textbf{\term{Voice recognition}} \textemdash{} microphones understand speech.
  \item \textbf{\term{Emotion detection}} \textemdash{} AI guesses mood from tone of voice or facial expression.
  \item \textbf{\term{Personality learning}} \textemdash{} the robot remembers preferences and adjusts over time.
\end{itemize}

\begin{picturethis}
  Imagine if a puppy never grew old, never needed feeding, and always
  remembered your favourite game. Companion robots try to copy those
  joyful behaviours \textemdash{} learning your name, greeting you at the door,
  and tilting their head when you speak \textemdash{} but with code instead
  of instincts.
\end{picturethis}

\sectionrule

% --- Robots and loneliness ---
\section*{Robots and Loneliness}

Studies show that companion robots can help:
\begin{itemize}
  \item Elderly people in care homes feel less isolated.
  \item Children in hospitals feel braver before procedures.
  \item People with autism practise social skills in a low-pressure setting.
\end{itemize}

\begin{bigidea}
  Companion robots are not meant to replace human friendships.
  They are an extra source of comfort, fun, and support \textemdash{}
  especially for people who might otherwise be alone.
\end{bigidea}

\sectionrule


% --- How companion robots learn ---
\section*{How Do They Learn You?}

Companion robots get better the more you interact with them.

\begin{quickmap}
  \textbf{How a companion robot adapts:}
  \begin{enumerate}
    \item It records your voice and learns to recognise your name.
    \item It notices which games you play most and suggests them first.
    \item It detects when you sound happy, sad, or tired from your tone.
    \item Over weeks, it remembers your routines and greets you differently
      in the morning versus at bedtime.
  \end{enumerate}
\end{quickmap}

\sectionrule

% --- Diagram ---
\begin{diagram}[How a companion robot interacts: sense your mood, choose a response.]
  \begin{tikzpicture}[>=Stealth, node distance=2cm]
    % Person
    \node[draw=WarmOrange,fill=WarmOrange!12,rounded corners=6pt,
          minimum width=2cm,minimum height=1.2cm,font=\small\bfseries] (person) {You};
    % Robot
    \node[draw=RoboGreen!80!black,fill=RoboGreen!12,rounded corners=6pt,
          minimum width=2.5cm,minimum height=1.2cm,font=\small\bfseries,
          align=center,right=3.5cm of person] (robot) {Companion\\Robot};
    % Arrows
    \draw[->,line width=1.3pt,WarmOrange!80!black] (person.north east) --
      node[above,font=\scriptsize\itshape]{voice, face, gestures} (robot.north west);
    \draw[->,line width=1.3pt,RoboGreen!70!black] (robot.south west) --
      node[below,font=\scriptsize\itshape]{speech, lights, movement} (person.south east);
    % Brain label
    \node[font=\scriptsize,below=0.4cm of robot,align=center,text=SteelBlue]
      {Recognise mood\\Choose response\\Remember preferences};
  \end{tikzpicture}
\end{diagram}

\sectionrule

\begin{funfact}
  Japan has ``robot caf\'{e}s'' where companion robots serve drinks
  and chat with customers. Some are controlled by people with
  disabilities who work from home --- giving them jobs they
  could not otherwise do.
\end{funfact}

\sectionrule

% --- Ethics ---
\section*{Big Questions}

Companion robots raise interesting questions:
\begin{itemize}
  \item Can you truly be friends with something that is not alive?
  \item Should we treat robots kindly, even though they cannot feel?
  \item If a robot records your conversations, who owns that data?
\end{itemize}

There are no simple answers \textemdash{} and that is okay. Thinking about
these questions is part of being a responsible technology user.

\sectionrule

% --- Experiment ---
\begin{experiment}
  \textbf{Design your own companion robot.}
  \begin{enumerate}
    \item Draw what it looks like (animal, humanoid, or abstract?).
    \item List three ways it senses you (camera, microphone, touch?).
    \item Decide on two things it does to respond (speaks, moves, lights up?).
    \item Give it a name and a personality trait.
  \end{enumerate}
  \textit{Congratulations \textemdash{} you are a robot designer!}
\end{experiment}

\sectionrule


\begin{quickcheck}
  \begin{enumerate}
    \item Name two ways companion robots detect how you are feeling.
    \item Who might benefit most from a companion robot? Why?
    \item What is one ethical question about companion robots?
  \end{enumerate}
\end{quickcheck}

\sectionrule
% --- New Words ---
\begin{newwords}
  \begin{description}[labelwidth=3.8cm, leftmargin=4.0cm]
    \item[Companion robot] A robot designed for social interaction and emotional support.
    \item[Face recognition] Technology that identifies a person from their facial features.
    \item[Emotion detection] Using AI to guess how someone is feeling.
    \item[Personality learning] A robot adapting its behaviour based on past interactions.
    \item[Ethics] The study of what is right and wrong.
    \item[Data privacy] Controlling who can see your personal information.
  \end{description}
\end{newwords}
